\chapter{相关工作与技术基础}

本章按方法家族而不是按年份逐篇罗列相关工作。这样组织的目的有两点：一是让读者先看清本文试图补哪一类缺口，二是避免把“隐藏指令发现”“形式化语义”“工具链基础设施”等本质不同的问题混写成一条松散文献链。

\paragraph{隐藏指令发现与指令空间审计。}
\tool{sandsifter} 是这一方向最具代表性的早期工作之一，它在 x86 上通过页错误分析与搜索空间裁剪执行大规模指令审计，发现了隐藏指令、硬件异常和多类工具链 bug \citep{domas2017breakingx86}。这项工作的意义首先在于证明了“处理器并不是可信黑盒”这一命题可以通过系统化实验方法落地验证；其次，它也为后续工作树立了一个重要边界：在变长 ISA 上，搜索空间剪枝本身就是研究问题。本文不尝试复现 x86 的复杂剪枝体系，而是将重点转向固定字长或更容易统一构造测试页的 RISC 风格 ISA，并把贡献放在执行控制、证据组织和可扩展接入上。

\paragraph{RISC 用户态扫描与差异 triage。}
\tool{iScanU} 是与本文最接近的先行工作。它把 \ptracebackend\ 与 \memcage\ 组织成 RISC 处理器上的用户态扫描框架，并通过信号与反汇编联合分析来识别未文档化行为 \citep{Dofferhoff_2020}。\tool{armshaker} 则聚焦 Armv8-A，强调从用户视角观察到的“隐藏行为”未必意味着真实硬件存在 undocumented instruction，其中相当一部分现象需要在 CPU、模拟器、反汇编器与操作系统之间做细致 triage \citep{Strupe_2020}。本文继承这两项工作的核心观察，但进一步把多 worker 编排、统一 raw 协议、Layer A/B/C 证据结构以及配置驱动生成链路纳入同一工程叙事。

\paragraph{authoritative semantics、合规测试与 formal verification。}
另一条研究路线并不直接做黑盒审计，而是试图提供权威语义或在其之上构建形式化验证。Sail 系列工作给出了 ARMv8-A、RISC-V 与 CHERI-MIPS 等 ISA 的大规模可执行语义模型，这些模型已经完整到能够支持真实系统软件启动 \citep{Armstrong_2019}。在此基础上，\tool{Islaris} 进一步展示了如何在 authoritative ISA semantics 之上开展机器码验证 \citep{Sammler_2022}。这些工作回答的是“在已有权威语义的前提下，如何验证程序行为”，而本文回答的是“当手头只有成品 CPU、现成工具链和用户态运行环境时，如何做第三方差异审计”。因此，两者并不竞争同一问题。它们对本文最重要的启发，是迫使我们诚实界定 ground truth：\capstone\ 在本文中是 practical decoder dependency，而不是 authoritative truth。

\paragraph{执行环境与工具链基础设施。}
QEMU 等模拟器和现代反汇编框架使大规模可复现实验成为可能。\citeauthor{bellard2005qemu} 的 QEMU 工作说明了跨 ISA 动态翻译和用户态模拟的工程基础 \citep{bellard2005qemu}。但本文并不把 QEMU 当作验证 oracle，而把它放在 Layer A 中作为受控开发与对照环境。这样的定位非常重要，因为本文在实验中已经观察到，模拟器、解码器和真实硬件之间可能形成多源差异。如果把 QEMU 或单一 decoder 直接写成“真值”，论文会陷入循环论证。基于同样原因，本文在结果解释中强调第二解码器、人工 triage 和后续实机补证，而不是把 Layer A 结果直接上升为硬件事实。

\paragraph{本文的位置。}
综合来看，本文相对已有工作的定位可以概括为三点。第一，相比 \tool{sandsifter}，本文不追求变长 ISA 的深度搜索裁剪，而强调面向 RISC 风格 ISA 的统一执行控制与扩展路线。第二，相比 \tool{iScanU}，本文把“多 ISA 工程化落地”“配置驱动生成链路”和“Layer A/B/C 分层证据结构”显式纳入主要贡献。第三，相比 Sail、\tool{Islaris} 等 formal 路线，本文的目标是面向现实硬件和现成工具链的黑盒审计，而不是证明实现满足形式化规范。

\placeholdertable{tab:related-compare}{相关工作与本文的对比占位}{[表 2-1 待补]\\建议列：工作、目标 ISA、执行控制、ground truth 依赖、是否包含实机证据、是否支持新架构低成本接入。\\该表将依据 \texttt{LITERATURE\_REVIEW.md} 在终稿中补齐。}
